
\section{Being and Identity}
\label{sec:being-and-identity}

\subsection{The Ontological Inseparability of Being and Identity}

The relation between being and identity constitutes one of the
foundational problems of ontology. Any account of what it means for
something to be must also address what makes that existent the existent
that it is. Being cannot be adequately considered without addressing
identity, just as identity cannot be understood independently of the
being whose identity it is.

The position developed here begins with a fundamental proposition:
being and identity are ontologically inseparable. They do not constitute
two independently established realities that subsequently enter into a
relationship. Nor does identity arise as an additional property
attributed to an entity that already exists in a fully constituted
ontological condition.

To say that something exists is, in this framework, to acknowledge
that it is determinate as the existent that it is. Its identity is not
an external qualification imposed upon its existence. It belongs to
the ontological actuality of that existence.

This proposition does not mean that the concepts of being and identity
are synonymous. Their conceptual distinction remains necessary.
Being concerns the actuality of existence, whereas identity concerns
the ontological self-determination through which an existent is the
existent that it is. These formulations identify analytically
distinguishable dimensions of an inseparable ontological actuality.

The distinction is therefore analytical rather than ontologically
separative. It allows the relation between existence and its
determinate character to be examined without converting them into
independent ontological components.

\subsection{Against the Priority of an Indeterminate Existent}

A possible account of existence assumes that an entity first exists
and subsequently acquires an identity through the addition of
properties, characteristics, or determinations. According to this
account, existence is initially considered independently of identity,
while identity is treated as a secondary condition that makes an
already existing entity identifiable or distinguishable.

The present ontology does not adopt this position.

An existent entirely devoid of identity would not constitute a
meaningful ontological starting point. If nothing could be said about
what makes an existent the existent that it is, then the supposed
existent would have been abstracted from the very condition through
which its existence is determinate.

This does not imply that every existent must possess a fully
articulated, externally accessible, or exhaustively describable
configuration. It means that the absence of a complete description
cannot be equated with the absence of identity.

The distinction is important. Ontological determination does not
depend upon whether an entity can be comprehensively characterized.
The limitations of a description do not establish a corresponding
indeterminacy in the existent itself.

Consequently, the distinction between an existent and its description
must not be confused with a distinction between existence and identity.
The former is legitimate; the latter, when interpreted as an
ontological separation, is rejected here.

\subsection{Being as Ontological Actuality}

Being is approached in this ontology as the actuality of existence,
rather than as a descriptive category, a linguistic designation, or
an abstraction constructed through conceptual operations.

An existent is not made real by being named, represented, classified,
or recognized. Its ontological actuality is not dependent upon the
presence of an entity capable of apprehending it.

This independence is essential to the scope of the present
investigation. The ontology of viability must establish its account
of existence without presupposing a knowing subject or making
existence conditional upon epistemic accessibility.

The fact that an existent may be unknown, incompletely described,
or inaccessible within a particular domain does not by itself
diminish its ontological actuality. Such limitations concern the
conditions under which an existent may become accessible within a
particular relation or domain. They do not constitute the ground of
its being.

Being must therefore be distinguished from the ways in which an
existent may be expressed, manifested, differentiated, or described.
These relations may be ontologically significant, but none can be
identified with existence as such without further argument.

The present section does not attempt to provide an exhaustive
definition of being. Its immediate purpose is more restricted:
to establish that being cannot be treated as an identity-free
ontological substrate.

\subsection{Identity as the Determinate Character of Being}

If being and identity are inseparable, the question arises as to
what identity contributes to the ontological constitution of an
existent.

Identity is not understood here as a list of properties, a collection
of observable features, or a descriptive profile. Such accounts may
characterize aspects of an existent, but they do not necessarily
establish what makes that existent the existent that it is.

Identity concerns the determinate character of existence itself.
It expresses the fact that an existent is not merely an unspecified
possibility but is ontologically constituted as something determinate.

This formulation must not be interpreted as reducing identity to
determination in its most general sense. Not every determination
necessarily constitutes identity, and not every feature associated
with an existent has the same ontological significance.

The distinction between constitutive determination and other forms
of determination requires a separate analysis. At this stage, the
essential claim is that identity cannot be detached from the
determinate actuality of being.

Identity is therefore neither an additional object within an entity
nor a separate entity standing alongside it. It is an ontological
dimension of the existent's own actuality.

\subsection{The Distinction Between Ontological Inseparability and Conceptual Identity}

The inseparability of being and identity does not eliminate the
possibility of distinguishing them conceptually.

Several ontological concepts may refer to dimensions of the same
existent without being interchangeable. The existence of an entity,
its identity, its unity, its individuation, and its particular
determinations may be examined separately while remaining related
within the same ontological actuality.

The distinction between conceptual differentiation and ontological
separation is therefore indispensable.

Conceptual differentiation makes it possible to investigate
different questions:

\begin{itemize}
    \item Being concerns the actuality of an existent.
    \item Identity concerns what makes that existent the existent
    that it is.
    \item Determination concerns the ways in which an existent is
    constituted as determinate.
    \item Unity concerns the ontological coherence of an existent.
    \item Individuation concerns its constitution as a distinct
    existent.
\end{itemize}

These questions are related, but their answers cannot be presumed
to be identical. In particular, the inseparability of being and
identity does not establish that identity is equivalent to unity,
individuation, or any particular form of determination.

The analytical distinction also prevents the opposite error:
treating every conceptually distinguishable dimension as an
independent ontological component.

The ontology developed in this work therefore distinguishes
without fragmenting and relates without collapsing.

\subsection{Identity and the Problem of an Underlying Substrate}

A further difficulty emerges when being is conceived as an
indeterminate substrate that remains ontologically prior to all
determinations.

Under this interpretation, identity may appear to be something
that happens to an underlying existent, rather than something
inseparable from the actuality of that existent.

The difficulty is not merely terminological. It concerns the
ontological status attributed to the supposedly underlying substrate.

If the substrate exists independently of every determination,
then the account must explain what makes it an existent rather
than an abstract possibility. If it possesses no identity whatsoever,
its status as a particular existent remains unexplained.

This ontology does not deny that abstraction may distinguish
existence from particular determinations for analytical purposes.
It rejects the further conclusion that such abstraction reveals
an independently existing, identity-free reality.

An analytical abstraction cannot, by itself, establish the
ontological independence of what it abstracts.

Accordingly, the concept of an underlying substrate must not be
used to reintroduce an ontological priority of being over identity
without an independent argument establishing that priority.

The present position is not that every question concerning
ontological grounds has already been resolved. Rather, it establishes
a constraint on subsequent explanations: no account may assume
an identity-free existent as its starting point without explaining
how such an existent could be ontologically actual.

\subsection{Being, Identity, and the Possibility of Difference}

The inseparability of being and identity also has implications
for the problem of difference.

If an existent is ontologically determinate, it cannot be considered
in complete abstraction from the conditions through which it is
distinct from other existents or from other possible configurations
of existence.

Nevertheless, identity must not be reduced to difference.
Difference concerns a relation or distinction, whereas identity
concerns the determinate actuality of the existent itself.

An entity may be distinguished from another without its identity
being exhausted by that distinction. Likewise, the fact that
different entities share certain determinations does not establish
that they possess the same identity.

The ontological relation between identity and difference therefore
requires an account that preserves both determination and
distinction without making either one reducible to the other.

This issue will be examined more extensively in the subsequent
analysis of identity and distinction.

\subsection{The Independence of Being and Identity from the Knowing Subject}

The ontological relation between being and identity must be
established independently of the conditions of knowledge.

An existent does not acquire being through recognition, nor does
its identity originate in the act of identifying it. The existence
of a knowing subject may introduce particular relations to an
existent, but such relations do not constitute the necessary
foundation of the existent's being or identity.

This position excludes an epistemic criterion from the foundational
definition of identity. An entity does not need to be known,
recognized, represented, or successfully distinguished by a
subject in order to possess identity.

It also excludes the assumption that the limits of knowledge
determine the limits of being. An existent may exceed any particular
description without thereby becoming ontologically indeterminate.

The present ontology consequently treats the knowing subject
as unnecessary to the establishment of the being--identity relation.
The analysis of how a knowing entity can apprehend existence,
including the conditions and limitations of such apprehension,
belongs to a distinct epistemological investigation.

This methodological boundary is not a denial of the reality of
knowledge. It is a distinction between the ontological constitution
of what exists and the epistemological conditions under which
existence may become known.

\subsection{Ontological Consequences of the Being--Identity Relation}

The inseparability of being and identity establishes several
constraints for the ontology of viability.

First, existence cannot be defined as an identity-free actuality
to which identity is subsequently added.

Second, identity cannot be treated as an external property,
descriptive label, or epistemic construction whose presence
depends upon recognition.

Third, conceptual distinctions between being, identity,
determination, unity, individuation, and difference must be
preserved without assuming that each distinction corresponds
to an independently existing ontological component.

Fourth, the incompleteness of a description must not be confused
with the absence or incompleteness of the existent's ontological
identity.

Fifth, no particular descriptive configuration, manifestation,
or invariant can be assumed to exhaust identity merely because
it provides a means of distinguishing or characterizing an existent.

These consequences do not yet provide a complete theory of
identity. They establish the ontological conditions within which
such a theory must be developed.

In particular, the relation between identity and constitutive
determination, the relation between identity and unity, and the
possibility of multiple identities associated with a single existent
remain questions requiring further examination.

\subsection{Concluding Formulation}

Being and identity are neither two independent ontological
realities nor two interchangeable concepts. They are analytically
distinguishable dimensions of an inseparable ontological actuality.

Being refers to the actuality of existence; identity refers to
the determinate character through which an existent is the existent
that it is. Neither has ontological priority over the other.
Identity is not subsequently added to being, and being is not
an identity-free substrate awaiting determination.

This position establishes a foundational principle for the
ontology of viability:

\begin{quote}
An existent does not first exist and then acquire an identity.
Its existence and its being-the-existent-that-it-is are
ontologically inseparable, although they remain conceptually
distinguishable.
\end{quote}

This principle does not resolve every question concerning the
constitution, unity, individuation, persistence, or multiplicity
of identity. It establishes the ontological ground from which
those questions must be investigated.

The subsequent analysis will therefore examine how identity
relates to determination, without presuming that the general
inseparability of being and identity has already settled the
specific structure of their constitutive relation.
